\documentclass[10pt, aspectratio=169]{beamer}

% Load the custom theme
\usetheme{violet}

% Title page information
\title[Short Title]{Your Main Presentation Title}
\subtitle{An Optional Subtitle for the Presentation}
\author{Your Name}
\institute{Your Institute}
\date{\today}

\begin{document}

% --- Title Page ---
\begin{frame}
    \titlepage
\end{frame}

% --- Section 1: Introduction ---
\section{Introduction}
\begin{frame}{Background and Motivation}
    \begin{itemize}
        \item First primary point introducing the topic and establishing relevance.
        \item Second point providing necessary context or literature review.
        \item Third point stating the core research question or objective.
    \end{itemize}
    
    \vspace{1em}
    \begin{block}{Key Concept}
        This block environment demonstrates the custom colors defined in the theme, utilizing the primary and secondary violet shades for emphasis.
    \end{block}
\end{frame}

% --- Section 2: Methodology ---
\section{Methodology}
\begin{frame}{Empirical Model}
    We estimate the following baseline model:
    
    \vspace{1em}
    $$
    Y_{it} = \alpha + \beta X_{it} + \gamma Z_{it} + \mu_i + \delta_t + \varepsilon_{it}
    $$
    \vspace{1em}
    
    \begin{itemize}
        \item $Y_{it}$: The dependent variable of interest.
        \item $X_{it}$: The primary explanatory variable.
        \item $Z_{it}$: A vector of time-varying control variables.
        \item $\mu_i, \delta_t$: Individual and time fixed effects.
    \end{itemize}
\end{frame}

% --- Section 3: Results ---
\section{Results}
\begin{frame}{Main Regression Output}
    \begin{table}
        \centering
        \begin{threeparttable}
            \caption{Estimated Effects on Primary Outcome}
            \begin{adjustbox}{max width=0.8\textwidth}
            \begin{tabular}{l c c}
                \toprule
                & (1) Baseline & (2) With Controls \\
                \midrule
                Variable of Interest ($X$) & 0.452*** & 0.389** \\
                                         & (0.112)  & (0.154) \\
                Control Variable ($Z$)     &          & 0.085 \\
                                         &          & (0.062) \\
                \midrule
                Observations               & 5,000    & 5,000   \\
                Fixed Effects              & Yes      & Yes     \\
                \bottomrule
            \end{tabular}
            \end{adjustbox}
            \begin{tablenotes}\footnotesize
                \item \textit{Notes:} Standard errors clustered at the group level in parentheses. *** p$<$0.01, ** p$<$0.05, * p$<$0.1.
            \end{tablenotes}
        \end{threeparttable}
    \end{table}
\end{frame}

% --- Section 4: Implementation ---
\section{Implementation}
% Note: The [fragile] option is REQUIRED for any frame containing verbatim code blocks
\begin{frame}[fragile]{Code Snippet Example}
    The following snippet demonstrates the monospace font rendering for code environments:
    
    \vspace{1em}
    \begin{block}{Pseudocode / Syntax}
\begin{verbatim}
import pandas as pd
import statsmodels.api as sm

# Fit a standard OLS model with clustered standard errors
model = sm.OLS(Y, sm.add_constant(X))
results = model.fit(cov_type='cluster', cov_kwds={'groups': df['id']})
print(results.summary())
\end{verbatim}
    \end{block}
\end{frame}

% --- Section 5: Conclusion ---
\section{Conclusion}
\begin{frame}{Summary}
    \begin{itemize}
        \item \textbf{Finding 1:} A brief summary of the first main result and its implications.
        \item \textbf{Finding 2:} A brief summary of the second main result.
        \item \textbf{Future Work:} Outline potential extensions or robustness checks for the current methodology.
    \end{itemize}
\end{frame}

\end{document}